% Strategic Management – Midterm Exam Phrases (one-pager)
% Compile: pdflatex midterm-exam-phrases.tex

\documentclass[8pt,landscape,a4paper]{article}
\usepackage[utf8]{inputenc}
\usepackage[T1]{fontenc}
\usepackage[margin=0.8cm]{geometry}
\usepackage{multicol}
\usepackage{enumitem}
\setlist{nosep,leftmargin=*}
\usepackage{xcolor}
\usepackage{microtype}
\pagestyle{empty}
\definecolor{eublue}{RGB}{0,51,153}

\begin{document}
\small
\begin{center}
{\large\bfseries\color{eublue} Strategic Management -- Midterm Exam Phrases}\\[0.2em]
\footnotesize Use these sentence patterns in answers. Name the framework, give reason, end with \emph{therefore}.
\end{center}
\vspace{0.3em}

\begin{multicols}{2}

\section*{\color{eublue}1. Master pattern}
\begin{itemize}
\item \emph{From a \textbf{[framework]} perspective, \textbf{[force/concept]} is \textbf{high/low} because \textbf{[reason]}, which makes the industry/situation \textbf{[attractive/unattractive]} for \textbf{[who]}, therefore \textbf{[strategic implication]}.}
\end{itemize}

\section*{\color{eublue}2. Internal vs external}
At the \textbf{external} level (PESTEL, Five Forces) we analyse \textbf{industry structure and profit potential}; at the \textbf{internal} level (RBV, VRIO) we explain \textbf{firm-level performance differences}. A good diagnosis combines both: the external view to understand the playing field, and the internal view to see whether the firm has the resources and capabilities to build and sustain competitive advantage on that field.

\section*{\color{eublue}3. CA in an unattractive industry}
An industry can be structurally unattractive (strong rivalry, powerful suppliers and substitutes), yet a firm may still obtain above-average returns if it controls \textbf{VRIO resources} (brand, technology, network effects) that rivals cannot easily replicate. In that case a negative industry effect is more than offset by a positive firm effect; exit is not always necessary if the firm occupies a protected niche.

\section*{\color{eublue}4. Intangibles \& accounts}
Conventional financial statements \textbf{understate firm value} because they record tangible assets at historical cost while most of the competitive edge lies in \textbf{intangibles} (human capital, reputation, culture, routines). These assets are decisive for VRIO yet are rarely fully capitalised, so the ``missing'' value is embedded in resources the accounts do not see.

\section*{\color{eublue}5. Grant's four (case template)}
For any case: (1) goals (long-term economic, social, reputational objectives); (2) understanding of the environment (industry structure, macro trends, culture); (3) resources and capabilities evaluated under VRIO; (4) implementation (are choices coherent with goals, environment, resource base?). Where these four elements are aligned, strategy is coherent and sustainable; where they diverge, performance suffers.

\section*{\color{eublue}6. Airlines \& industry structure}
Given capital intensity, regulation and exposure to powerful suppliers and substitutes, airlines operate in a high-risk, low-margin strategic context. Many cost drivers are industry-wide (fuel, aircraft, regulation), so sustainable advantage tends to arise from \textbf{network design, hub control, brand and loyalty programmes} rather than from input prices alone.

\section*{\color{eublue}7. Governance, stakeholders, agency}
The separation between ownership and control creates an \textbf{agency problem}: managers may pursue private benefits (power, pay, empire building) that do not maximise shareholder wealth, so incentive schemes and board oversight are needed to align interests. A pure shareholder view focuses on return on equity and cost of capital, whereas a stakeholder view broadens performance to include social and environmental externalities; this affects which strategies are seen as legitimate.

\section*{\color{eublue}8. SWOT done right}
A useful SWOT does not list generic strengths and opportunities; it identifies \textbf{specific internal attributes} that plausibly explain why customers choose the firm and matches them with \textbf{concrete external trends}. S/W are internal (resources, capabilities); O/T are external (PESTEL, Five Forces). Each point should be linked to competitive advantage, not just description.

\section*{\color{eublue}9. Clusters \& industry effect}
Clusters show how location can modify industry economics: geographical concentration of firms, suppliers and universities can reduce transaction costs, accelerate innovation and soften certain Five Forces, increasing profit potential. Observed profitability is thus a combination of \textbf{industry effect} (how attractive the sector is) and \textbf{firm effect} (strategy and resources).

\section*{\color{eublue}10. Performance ``compared to what?''}
Performance makes sense only \textbf{relative to a referent}: industry average, main rival, the firm's own past, or its stated intentions. A raw ROA or ROE number is not ``good'' or ``bad'' in isolation; it becomes meaningful when compared to these benchmarks.

\section*{\color{eublue}11. Cost of equity \& opportunity cost}
For listed firms, shareholders require a minimum return equal to the \textbf{cost of equity}: $R_f + \beta(\text{Market} - R_f)$. Satisfaction occurs when realised return is at least equal to this opportunity cost or beats the industry average; otherwise capital would have been better deployed elsewhere.

\section*{\color{eublue}12. Porter's generic strategies}
According to Porter, firms can seek advantage by \textbf{low-cost}, \textbf{differentiation}, or \textbf{focus}. Low-cost wins with price-sensitive customers; differentiation creates inelastic demand through brand, quality or service; focus serves a narrow segment. Being ``stuck in the middle'' (neither clearly low-cost nor differentiated) tends to produce low margins and weak loyalty.

\section*{\color{eublue}13. Barriers to entry (natural vs artificial)}
Barriers to entry can be \textbf{natural} (economies of scale, know-how, learning curves, capital intensity) or \textbf{artificial} (licences, regulation, tariffs, patents). High entry barriers protect incumbents and support sustained profit, while low barriers invite new competitors and erode margins.

\section*{\color{eublue}14. Buyer and supplier power}
Buyer power is high when few large buyers, low switching costs and full information allow customers to push prices down; supplier power is high when inputs are concentrated, differentiated and costly to replace, allowing suppliers to push input prices up. In both cases, high power on the other side of the transaction reduces the focal firm's share of value.

\section*{\color{eublue}15. Stock openers}
\begin{itemize}
\item \emph{``This question is essentially about \textbf{[topic]} (e.g.\ industry attractiveness / internal vs external analysis / CA in an unattractive industry).''}
\item \emph{``From a \textbf{Porter Five Forces} perspective\ldots''} \quad \emph{``Using the \textbf{resource-based view (RBV)}\ldots''}
\item \emph{``At the \textbf{macro} level (PESTEL)\ldots''} \quad \emph{``At the \textbf{firm} level (VRIO)\ldots''}
\end{itemize}

\section*{\color{eublue}16. Connective phrases}
\begin{itemize}
\item \emph{``First\ldots Second\ldots Third\ldots''} (structure your points clearly).
\item \emph{``Moreover\ldots''} \quad \emph{``However\ldots''} \quad \emph{``In addition\ldots''}
\item \emph{``In economic terms, this implies that\ldots''} \quad \emph{``From a strategic management point of view\ldots''}
\item \emph{``In summary, the combined effect of these forces is that\ldots''}
\end{itemize}

\end{multicols}

\end{document}

